\section{模型评价与推广}

\subsection{模型优点}

（1）模型层次递进，逻辑严密。从问题一的基础功率平衡计算，到问题二的离散优化，再到问题三的连续优化和问题四的离网储能配置，每个子问题都在前一个的基础上放松约束或增加维度，形成了完整的研究体系。各子问题之间的递进关系（离散$\subset$连续、联网$\supset$离网）保证了结果的可比性和一致性。

（2）求解方法与问题特征匹配。问题一采用代数计算（无优化需求），问题二利用目标函数可分离性证明贪心最优，问题三建立标准LP利用成熟求解器，问题四采用两层优化（外层搜索+内层LP）。每种方法都有明确的数学依据，避免了"杀鸡用牛刀"的过度建模。

（3）物理约束建模完整。功率平衡方程、购售电互斥条件、设备功率上下限、储能状态方程和周期性约束等物理约束均在模型中严格体现，计算结果通过了能量守恒和边界条件验证，确保了物理合理性。

（4）结果分析维度丰富。不仅给出了最优解的数值，还从成本分解、绿电指标、设备利用率、全年分布特征等多个维度进行了深入分析，揭示了各因素之间的内在联系和权衡关系。

\subsection{模型局限性}

（1）时间分辨率较粗。采用1小时时段假设（分段恒定功率），未考虑时段内的功率波动。实际运行中风电和光伏的分钟级波动可能对设备调节提出更高要求，建议在精细化研究中采用15分钟或5分钟时段。

（2）设备启停约束简化。实际工程中电解槽和合成氨装置的启停存在最小运行时间、最小停机时间、爬坡速率等约束，本模型中未予考虑。引入这些约束将使问题成为混合整数规划（MIP），计算复杂度显著增加。

（3）氢气储存环节缺失。本模型假设制氢和合成氨同步运行（即产即用），未考虑中间氢气缓冲储罐的解耦作用。引入氢储环节可以允许制氢和制氨独立调度，增加优化自由度。

\subsection{模型推广}

本模型框架可推广至以下方向：（1）多园区集群协调调度，引入园区间的功率互济和虚拟电厂聚合；（2）考虑碳交易市场的经济性分析，将碳排放权价格纳入成本函数；（3）长时储能（如压缩空气、液流电池）配置优化，突破锂电池4小时容量限制；（4）风光功率预测不确定性下的鲁棒优化或随机规划，提升调度方案的抗风险能力。
